% ===================================================================
%  Dodatek K: Asymptotyczny ogon solitonu — analityczne skalowanie
%             A_tail i program asymptotyki dopasowanej
%  TGP v1 — Sesja v34 (2026-03-28)
%  Adresuje otwarte problemy: O-J1, O-J2, O-J3
% ===================================================================
\section{Asymptotyczny ogon solitonu: skalowanie, WKB i predykcja tau}%
\label{app:wkb-atail}

Niniejszy dodatek uzupełnia Dodatek~J (\S\ref{app:ogon-masy})
o~trzy elementy formalne:
\begin{enumerate}[nosep]
  \item \textbf{Ścisłe wyprowadzenie formy ogona} $g(r)-1 \sim (B\cos r + C\sin r)/r$
        z~linearyzacji ODE solitonu (zamykając O-J1 częściowo).
  \item \textbf{Argument skalowania $A_{\rm tail} \propto (g_0-g^*)^4$}
        z~charakterystyki potencjału $V(g) = g^3/3-g^4/4$.
  \item \textbf{Predykcja masy leptonu tau} $g_0^\tau$ i~stosunku
        $r_{31} = m_\tau/m_e$ (status O-J3).
\end{enumerate}

Duch teorii zostaje zachowany: wszystkie wyprowadzenia korzystają wyłącznie
z~aksjomatu $N0$-2 ($\alpha=2$, operator $\mathcal{D}$) i~warunku próżniowego
$N0$-5 ($\beta=\gamma$, $V(g) = g^3/3 - g^4/4$).

\medskip
\noindent\fbox{\parbox{\linewidth}{%
\textbf{Nota kanoniczna (2026-04-07).}
Niniejszy dodatek u\.zywa Formulacji~B ($V = g^3/3-g^4/4$, $f(g) = 1+4\ln g$, $g_0^e \approx 1{,}24$, punkt duchowy $g^* = e^{-1/4}$).
W~sformu\l{}owaniu kanonicznym: $K(g) = g^4$, $P(g) = (\beta/7)g^7 - (\gamma/8)g^8$,
$g_0^e = 0{,}86941$, brak punktu duchowego.
Algebra Koide'a i~r\'ownania master F1--F12 s\k{a} niezale\.zne od tego wyboru.}}
\medskip

% -------------------------------------------------------------------
\subsection{Linearyzacja ODE przy próżni: forma ogona}%
\label{app:K-linearyzacja}
% -------------------------------------------------------------------

\subsubsection{Równanie solitonu TGP}

ODE solitonu sferycznie symetrycznego (por.\ \S\ref{app:J-profil}):
\begin{equation}\label{eq:K-ode}
  f(g)\,g'' + \frac{2}{r}\,g' = V'(g),
  \qquad
  f(g) = 1 + 2\alpha\ln g,\quad \alpha = 2,
  \qquad
  V'(g) = g^2(1-g),
\end{equation}
z~warunkami: $g(0) = g_0$, $g'(0) = 0$, $g(r)\to 1$ dla $r\to\infty$.

\subsubsection{Punkt osobliwy kinetyczny $g^*$}

Czynnik $f(g)$ znika przy
\begin{equation}\label{eq:K-gstar}
  g^* = \exp\!\bigl(-\tfrac{1}{2\alpha}\bigr) = e^{-1/4} \approx 0{,}77880,
\end{equation}
co tworzy \textbf{barierę kinetyczną}: przy $g \to g^*$ efektywne
przyspieszenie $g'' = V'(g)/f(g) \to \infty$, a~trajektoria nie może
przekroczyć $g^*$ (twierdzenie~\ref{thm:N0-duality} i~uw.~\ref{rem:J-Atail-fiz}).
Przestrzeń rozwiązań jest zatem $g(r)\colon [0,\infty)\to (g^*,\infty)$.

\subsubsection{Linearyzacja przy próżni $g = 1$}

Dla dużych $r$, gdy $g\to 1$, piszemy $g = 1 + u$ z~małym $u$:
\begin{align}
  f(1+u) &= 1 + 2\alpha\ln(1+u) \approx 1 + 2\alpha u + O(u^2), \\
  V'(1+u) &= (1+u)^2(-u) = -u - 2u^2 - u^3 \approx -u + O(u^2).
\end{align}
Równanie~\eqref{eq:K-ode} do pierwszego rzędu w~$u$:
\begin{equation}\label{eq:K-linear}
  u'' + \frac{2}{r}\,u' + u = 0.
\end{equation}
Jest to \textbf{równanie sferyczne Bessela rzędu 0} z~jednostkową
częstością przestrzenną ($m_{\rm sp} = 1$ w~jednostkach znormalizowanych
przez $\gamma = 1$, por.\ stwierdzenie~\ref{prop:N0-6-from-N0-5}).

\begin{proposition}[Ogon solitonu TGP]\label{prop:K-tail-form}\label{def:A-tail}
Rozwiązanie r\'ownania~\eqref{eq:K-linear} regularnego w~$r=0$
i~zanikającego dla $r\to\infty$ ma postać:
\begin{equation}\label{eq:K-tail}
  u(r) \;=\; g(r) - 1
  \;\sim\; \frac{B\cos r + C\sin r}{r}
  \qquad (r \to \infty),
\end{equation}
gdzie $B,C$ s\k{a} sta\l{}ymi zale\.znymi od~$g_0$ wyznaczanymi
przez warunki pocz\k{a}tkowe w~$r=0$ (,,strzał'').
Amplituda ogona:
\begin{equation}\label{eq:K-Atail}
  A_{\rm tail}(g_0) = \sqrt{B(g_0)^2 + C(g_0)^2}, \qquad
  \phi_{\rm tail}(g_0) = \mathrm{atan2}(B,C).
\end{equation}
\end{proposition}

\begin{proof}
Podstawienie $w = r\cdot u$ przekształca~\eqref{eq:K-linear} w~$w'' + w = 0$,
którego ogólne rozwiązanie to $w = B\cos r + C\sin r$,
skąd $u = (B\cos r + C\sin r)/r$.
Warunek $u\to 0$ dla $r\to\infty$ jest automatycznie spełniony
(nie wymaga $B = C = 0$).
\end{proof}

\begin{remark}[Masa bozonu przestrzennego $m_{\rm sp}$]
\label{rem:K-msp}
W~pełnym modelu TGP z~nieliniowością $\beta\neq 0$
linearyzacja przy $g=1$ daje $V''(1) = -\beta$ (gdzie~$\beta = \gamma$
z~N0-5, uw.~\ref{prop:vacuum-condition}).
Skuteczna masa bozonu przestrzennego:
\begin{equation}\label{eq:K-msp-full}
  m_{\rm sp} = \sqrt{\beta} = \sqrt{\gamma}.
\end{equation}
W~jednostkach, w~kt\'orych $\gamma = 1$ (stosowanych w~dodatku~J),
$m_{\rm sp} = 1$ i~ogon oscyluje z~jednostkową d\l{}ugością fali.
\end{remark}

% -------------------------------------------------------------------
\subsection{Argument skalowania $A_{\rm tail} \propto (g_0-g^*)^4$}%
\label{app:K-skalowanie}
% -------------------------------------------------------------------

\subsubsection{Wynik numeryczny}

Skrypty \texttt{ex57} (sesja v33+) i~\texttt{ex62} (sesja v34) potwierdzają:
\begin{equation}\label{eq:K-power-law}
  A_{\rm tail}(g_0) \;\approx\; c_A\,(g_0 - g^*)^{4{,}12\pm 0{,}05},
  \qquad
  c_A \approx 0{,}35,
  \qquad
  g_0 \in [1{,}1,\, 2{,}5].
\end{equation}
Wynik jest \textbf{niezależny od~$\alpha$}: skrypt \texttt{ex62}~(M4)
stwierdza, że dla $\alpha \in [1{,}0;\, 3{,}0]$ wykładnik pozostaje
blisko~4.

\subsubsection{Argument heurystyczny: stopień potencjału}

\begin{proposition}[Wykładnik z~potencjału --- \textbf{obalony}]\label{prop:K-exponent}
W~heurystycznym przybliżeniu amplituda promieniowania solitonu
miałaby skalować się jak kwadrat energii wewnętrznej:
\begin{equation}\label{eq:K-heuristic}
  A_{\rm tail} \;\sim\; E_{\rm sol}(g_0)^2,
  \qquad
  E_{\rm sol}(g_0) \;\propto\; (g_0 - 1)^2
  \;\;\text{dla}\;\; g_0 \to 1^+,
\end{equation}
skąd $A_{\rm tail} \propto (g_0-1)^4$.

\medskip\noindent
\textbf{Status: OBALONY (ex150, 2026-04-05).}
Skrypt \texttt{ex150} testuje argument numerycznie dla potencjałów
$V_n(g) = g^3/3 - g^n/n$ ($n = 3, 4, 5, 6$):
\begin{itemize}[nosep]
  \item Fit $A_{\rm tail}$ vs $(g_0-1)$: daje $k \approx 1{,}0$
    dla \emph{każdego}~$n$ (po prostu linearyzacja).
  \item Fit $A_{\rm tail}$ vs $(g_0-g^*)$: daje $k \approx 1{,}8$
    (nie~4), niezależnie od stopnia~$n$.
  \item Wykładnik NIE rośnie z~$n$ (nachylenie $dk/dn = 0{,}03$).
\end{itemize}
Skalowanie $A \propto (g_0-g^*)^{4{,}12}$ (eq.~\eqref{eq:K-power-law})
pochodzi ze skryptów \texttt{ex57}/\texttt{ex62}, które nie przetrwały
w~repozytorium --- wynik jest niereprodukowalny.
\end{proposition}

\begin{remark}[Mechanizm]
Energia solitonu pochodzi z~potencjału $V(g)$. Dla małych odchyleń od~$g=1$:
\begin{align}
  V(1+\delta) &= V(1) + V'(1)\delta + \tfrac{1}{2}V''(1)\delta^2 + \ldots
  = \tfrac{1}{12} - \tfrac{1}{2}\delta^2 + O(\delta^3),
\end{align}
gdzie $V(1) = 1/3 - 1/4 = 1/12$ i~$V''(1) = 2 - 3 = -1$
(minimum destabilizujące --- soliton jest solitonem niezaciągającym,
utrzymywanym przez barierę $g^*$).
Energia kinetyczna solitonu przy wejściu w~rejon próżni ($g\to 1$)
skaluje się jak $\delta^2 \sim (g_0-1)^2$ przez zachowanie energii
w~przybliżeniu jednorodnym.
Promieniowanie (amplituda ogona) jest proporcjonalne do kwadratu
tej energii: $A_{\rm tail} \sim (g_0-1)^4$.

Potencjał $V(g) = g^3/3 - g^4/4$ \textbf{ma stopień~4}
(dominujący człon $-g^4/4$). Jest to kluczowe: gdyby $V$ miał
stopień~6, oczekiwalibyśmy wykładnika~6; gdyby~2 --- wykładnik~2.
Wynik $n \approx 4$ jest wewnętrznie spójny z~N0-5 ($\beta = \gamma$
daje potencjał stopnia czwartego).
\end{remark}

\subsubsection{Całka WKB}

Alternatywne ujęcie przez całkę WKB:
\begin{equation}\label{eq:K-SWKB}
  S_{\rm WKB}(g_0)
  = \int_{g^*}^{g_0} \sqrt{\frac{V(g)}{|f(g)|}}\,dg.
\end{equation}
Przy $g \to g^*$: $f(g) \approx (2\alpha/g^*)(g-g^*)$, więc
podcałkowa $\sim 1/\sqrt{g-g^*}$ --- całkowalna.

\begin{result}[Skalowanie $S_{\rm WKB}$]%
\label{res:K-SWKB}
Skrypt~\texttt{ex62}~(M2) wyznacza numerycznie:
\begin{equation}\label{eq:K-SWKB-scaling}
  S_{\rm WKB}(g_0) \;\propto\; (g_0 - g^*)^{m_{\rm WKB}},
  \qquad
  m_{\rm WKB} \;\approx\; 1{,}5,
\end{equation}
co odpowiada oczekiwaniu geometrycznemu
$\int_0^\delta dx/\sqrt{x} \sim \sqrt{\delta}$ przy barierze
połączonemu z~działaniem potencjału $V(g^*) \sim g_0-g^*$:
$S_{\rm WKB} \sim \sqrt{(g_0-g^*)\cdot\mathrm{const}} \propto (g_0-g^*)^{3/2}$.
Natomiast $A_{\rm tail} \not\propto S_{\rm WKB}^p$ wprost ---
zależność jest złożoną kombinacją liniową i~nieliniową
(wynik M2 z~\texttt{ex62}).
\end{result}

\begin{problem}[O-J1: Analityczne $A_{\rm tail}(g_0)$]\label{prob:K-OJ1}
Rygorystyczne wyprowadzenie funkcji $A_{\rm tail}(g_0)$
przez metodę asymptotyki dopasowanej:
\begin{enumerate}[nosep]
  \item \textbf{Warstwa zewnętrzna}: rozwiązanie liniowe~\eqref{eq:K-tail}
        dla dużych $r$.
  \item \textbf{Warstwa wewnętrzna} przy $g \sim g^*$: ODE nieliniowe
        z~$f \approx f'_*(g-g^*)$, rozwiązanie przez zmienne Painlevé.
  \item \textbf{Dopasowanie}: połączenie asymptotyk daje $B(g_0)$ i~$C(g_0)$
        jako funkcje $g_0$ przez formułę połączeniową.
\end{enumerate}
Oczekiwany wynik: $A_{\rm tail} = c_A(g_0-g^*)^4[1+O((g_0-g^*))]$,
gdzie $c_A$ jest wyznaczone z~warunków dopasowania.
\end{problem}

% -------------------------------------------------------------------
\subsection{Warunek kwantowania $B_{\rm tail} = 0$ i~granice sektorów}%
\label{app:K-kwantowanie}
% -------------------------------------------------------------------

\begin{remark}[Interpretacja topologiczna sektorów]%
\label{rem:K-sektory}
Przestrzeń trajektorii $\mathcal{T} = \{g\colon [0,\infty) \to (g^*,\infty)\}$
dzieli się na klasy homotopii ze względu na liczbę $n_{\rm cross}$
przekroczeń poziomu $g = 1$ (z~góry na dół):
\begin{itemize}[nosep]
  \item $n_{\rm cross} = 0$: \textbf{sektor elektronu} --- trajektoria
    maleje z~$g_0$ do minimum $g_{\rm min} > 1$, po czym dąży do~$g=1$
    od~góry lub równoważnie wykonuje jedno półosiągnięcie.
  \item $n_{\rm cross} = 1$: \textbf{sektor mionu} --- trajektoria
    przekracza $g=1$ raz (z~góry), ma minimum $g_{\rm min} < 1$,
    po czym wraca do $g=1$ od~dołu.
  \item $n_{\rm cross} = k$: \textbf{sektor $k$-tego leptonu} ---
    $k$~przekroczeń poziomicy próżniowej.
\end{itemize}
Skrypt~\texttt{ex60}~(T3) identyfikuje wartości $g_0$ na~granicach sektorów
(patrz tablica w~wynikach).
\end{remark}

Warunek $B_{\rm tail}(g_0) = 0$ (hipoteza H1, stwierdzenie~\ref{prop:J-quantization})
wybiera wyróżnione wartości $g_0$ w~każdym sektorze:
soliton jest w~\textbf{rezonansie z~substratem} --- amplituda
fali wychodzącej (kosinusowej) zanika.

\begin{result}[Zera $B_{\rm tail}$ z~\texttt{ex60}]%
\label{res:K-zeros}
Skrypt~\texttt{ex60}~(T5) wyznacza kolejne zera $B_{\rm tail}(g_0)$
dla $g_0 \in [1{,}08;\, 4{,}0]$:
\begin{center}\footnotesize
\begin{tabular}{@{}ccccc@{}}
\toprule
$n$ & $g_0^{(n)}$ & Sektor & Identyfikacja & Status \\
\midrule
1 & $1{,}230$ & $n_{\rm cross}=0$ & elektron & stw.~\ref{prop:J-quantization} \\
2 & $\approx 2{,}08$ & granica & mion? & \statuslabel{Hipoteza} \\
3 & do wyznaczenia & $n_{\rm cross}=1$ & tau? & \statuslabel{Program} (O-J2/J3) \\
\bottomrule
\end{tabular}
\end{center}

\medskip\noindent
\textbf{Uwaga}: hipoteza złotej proporcji~\ref{hyp:J-golden-topo}
sugeruje $g_0^\mu = 2{,}00$, natomiast warunek $B_{\rm tail} = 0$
daje $g_0 \approx 2{,}08$ (odchylenie~4\%). Rozstrzygnięcie między
tymi dwoma zasadami selekcji pozostaje otwarte (\textbf{O-J2}).
Fizycznie obie zasady mogą być kompatybilne: $g_0^e = 1{,}24 \approx 1{,}230$
(1\%) i~$g_0^\mu = 2{,}00 \approx 2{,}08$ (4\%) --- obie z~dokładnością
rzędu~$1$--$4\%$, wpisującą się w~precyzję hipotezy masy~\ref{hyp:J-mass-Atail4}
(odchylenie 1{,}6\%).
\end{result}

% -------------------------------------------------------------------
\subsection{Predykcja masy tau: Ścieżka~9 (rozszerzenie)}%
\label{app:K-tau}
% -------------------------------------------------------------------

\subsubsection{Cel: wyznaczenie $g_0^\tau$}

Z~hipotezy masowej~\ref{hyp:J-mass-Atail4}: $m \propto A_{\rm tail}^4$,
stosunek $r_{31} = m_\tau/m_e = 3477{,}15$ (PDG~2024) wymaga:
\begin{equation}\label{eq:K-tau-condition}
  \left(\frac{A_{\rm tail}(g_0^\tau)}{A_{\rm tail}(g_0^e)}\right)^4
  = r_{31} = 3477{,}15.
\end{equation}

\subsubsection{Kandydaci}

\begin{center}\footnotesize\label{tab:K-tau-candidates}
\begin{tabular}{@{}llp{4.0cm}cc@{}}
\toprule
\textbf{Kandydat} & $g_0^\tau$ & \textbf{Zasada selekcji}
  & $(A_\tau/A_e)^4$ & $\Delta r_{31}$ [\%] \\
\midrule
C2: $g_0^\mu \cdot\phi$ & $2{,}00 \times 1{,}618 = 3{,}236$ &
  złota proporcja --- 2 kroki & do wyznaczenia & \\
C3: $g_0^e \cdot\phi^2$ & $1{,}24 \times 2{,}618 = 3{,}247$ &
  złota proporcja --- 2 potęgi & do wyznaczenia & \\
C4: Odwrotne & $\approx 2{,}34$ &
  numeryczne rozwiązanie~\eqref{eq:K-tau-condition} & $3477{,}15$ & $0\%$ \\
C5: Fit $A_{\rm tail}$ & $\approx 2{,}34$ &
  $c_A(g-g^*)^{4{,}12} = A_{\rm tail}^{\rm target}$ & $\approx r_{31}$ & $\lesssim 5\%$ \\
C1: $B_{\rm tail}=0$ & trzecie zero & warunek kwantowania & do wyznaczenia & \\
\bottomrule
\end{tabular}
\end{center}

\begin{remark}[Kandydat C4 jako predykcja bazowa]%
\label{rem:K-C4}
Kandydat~C4 (numeryczne odwrócenie~\eqref{eq:K-tau-condition})
daje \textbf{predykcję}: $g_0^\tau \approx 2{,}34 \pm 0{,}05$.
Stosunki:
\begin{equation}\label{eq:K-g0-ratios}
  \frac{g_0^\tau}{g_0^\mu} \approx \frac{2{,}34}{2{,}00} = 1{,}17,
  \qquad
  \frac{g_0^\tau}{g_0^e} \approx \frac{2{,}34}{1{,}24} = 1{,}89
  \approx \phi^{1{,}3}.
\end{equation}
Złota proporcja \textbf{nie} reprodukuje $r_{31}$ dla~$k=2$
(por.\ uwaga~\ref{rem:J-tau}): $(\phi^2)^{4{,}12} \gg 3477$.
Mechanizm wyznaczający $g_0^\tau$ różni się od mechanizmu dla~$g_0^\mu$.
\end{remark}

\begin{hypothesis}[Predykcja masy tau]\label{hyp:K-tau}
Masa leptonu tau wynika z~Ścieżki~9 przez:
\begin{equation}\label{eq:K-tau-prediction}
  m_\tau = m_e \times \left(\frac{A_{\rm tail}(g_0^\tau)}{A_{\rm tail}(g_0^e)}\right)^4
  \approx 3477 \times m_e \approx 1777\;\mathrm{MeV},
\end{equation}
gdzie $g_0^\tau \approx 2{,}34$ jest wyznaczone z~warunku odwrotnego~C4.
Odchylenie od PDG ($m_\tau = 1776{,}86$~MeV): $< 0{,}1\%$ przez konstrukcję.

Weryfikowalna zasada selekcji $g_0^\tau$ --- niezależna od~odwrotnej
kalkulacji~C4 --- pozostaje \textbf{otwartym problemem O-J3}.
Kandydaci to: (i)~trzecie zero~$B_{\rm tail}(g_0)$,
(ii)~warunek granicy sektora $n_{\rm cross}=1\to 2$,
(iii)~nowa zasada samodopasowania (program).

Status: \statuslabel{Hipoteza} ilościowa.
\end{hypothesis}

\subsubsection{Zgodność z formułą Koidego}

Dla mas elektron-mion-tau z~Ścieżki~9 możemy sprawdzić formułę Koidego:
\begin{equation}\label{eq:K-Koide-check}
  Q_K = \frac{(\sqrt{m_e} + \sqrt{m_\mu} + \sqrt{m_\tau})^2}
             {\tfrac{2}{3}(m_e + m_\mu + m_\tau)}
  \stackrel{?}{=} 1.
\end{equation}
Z~danymi PDG: $Q_K = 1{,}000000 \pm 10^{-6}$
(formuła Koidego dokładna na~poziomie $10^{-6}$).
Z~predykcją Ścieżki~9 ($m_\tau^{\rm TGP} \approx 1777$~MeV):
\begin{equation}
  |Q_K^{\rm TGP} - 1| \lesssim 10^{-4}
  \quad(\text{zależy od precyzji wyznaczenia }g_0^\tau).
\end{equation}
Zgodność z~Koidem dla TGP jest \textbf{testem spójności}:
predykcja $g_0^\tau$ powinna reprodukować nie tylko $r_{31}$,
ale i~relację Koidego do tej samej precyzji.
Skrypt~\texttt{ex61}~(sekcja~Bonus) weryfikuje tę zgodność.

\begin{remark}[v41: ODE substratowe --- predykcja \emph{wszystkich} mas leptonów]%
\label{rem:K-substrate-tau}
Skrypt \texttt{tau\_uv\_completion.py} (2026-03-30) rozwiązuje ODE
substratowe~\eqref{eq:J-substrate-ode} ($K_{\rm sub} = g^2$, ghost-free).
Wyniki:
\begin{center}\footnotesize
\begin{tabular}{@{}lcccc@{}}
\toprule
  Lepton & $g_0$ & $(A/A_e)^4$ & PDG & $\Delta$ \\
\midrule
  $e$ & $0{,}8694$ (z $\varphi$-FP) & $1$ & $1$ & --- \\
  $\mu$ & $1{,}4068 = \varphi\cdot g_0^*$ & $206{,}77$ & $206{,}77$ & $0{,}04\%$ \\
  $\tau$ & $1{,}7294$ & $3477{,}2$ & $3477{,}2$ & $0{,}001\%$ \\
\bottomrule
\end{tabular}
\end{center}
$m_\tau^{\rm pred} = 1776{,}83$~MeV (PDG: $1776{,}86$), $\Delta = -14$~ppm.
Formuła Koidego: $K = 0{,}66666$ ($\Delta = -11$~ppm od dokładnego $2/3$).

Stosunki w~ODE substratowym:
$g_0^\tau/g_0^e = 1{,}989$, $g_0^\tau/g_0^\mu = 1{,}229$, $g_0^\tau/(g_0^e\cdot\varphi^2) = 0{,}760$.
Żaden prosty stosunek z~$\varphi$ nie reprodukuje $g_0^\tau$ --- zasada
selekcji tauonu (O-J3) pozostaje \textbf{otwarta}.

\textbf{Uwaga}: ODE substratowe jest \emph{jedynym} wariantem ODE zdolnym
reprodukować $r_{31}$. Pełne ODE ($f(g) = 1+4\ln g$) obcina spektrum barierą
ducha przy $g_0^{\rm crit} \approx 1{,}63$, co ogranicza $\max(A/A_e)^4 \approx 1571$
--- zbyt mało o~czynnik $2{,}2\times$.
\end{remark}

% -------------------------------------------------------------------
\subsection{Status i~otwarte problemy (sesja v34)}%
\label{app:K-status}
% -------------------------------------------------------------------

\begin{center}\footnotesize
\begin{tabular}{@{}p{5.0cm}lp{4.0cm}p{2.2cm}@{}}
\toprule
\textbf{Zdanie} & \textbf{Status} & \textbf{Warunki / założenia} & \textbf{Odsyłacz} \\
\midrule

Ogon solitonu $\sim (B\cos r + C\sin r)/r$
  & \statuslabel{Twierdzenie}
  & linearyzacja ODE~\eqref{eq:K-ode} przy $g=1$; $m_{\rm sp}=1$
  & stw.~\ref{prop:K-tail-form} \\[3pt]

$A_{\rm tail}(g_0) \propto (g_0-g^*)^{4{,}12}$
  & \statuslabel{Propozycja}
  & numeryc.\ \texttt{ex57}, \texttt{ex62}; zakres $g_0 \in [1{,}1;\,2{,}5]$
  & stw.~\ref{prop:J-scaling} \\[3pt]

Wykładnik $n \approx 4$: stopień potencjału $V$
  & \statuslabel{Obalony}
  & \texttt{ex150}: $k \approx 1{,}0$ vs $(g_0-1)$; brak korelacji $k$--$n$
  & uw.~\ref{prop:K-exponent} \\[3pt]

$g_0^\tau \approx 2{,}34$ (predykcja Ścieżki~9)
  & \statuslabel{Hipoteza}
  & hipoteza~\ref{hyp:J-mass-Atail4} + odwrotne równanie~C4
  & hip.~\ref{hyp:K-tau} \\[3pt]

Zasada selekcji $g_0^\tau$ (niezależna od~C4)
  & \statuslabel{Program}
  & zera $B_{\rm tail}$, granice sektorów, samopod.
  & O-J3 \\[3pt]

Rygorystyczne wyprowadzenie $A_{\rm tail}(g_0)$
  & \statuslabel{Program}
  & asymptotyka dopasowana; warstwa przy $g^*$
  & O-J1, prob.~\ref{prob:K-OJ1} \\[3pt]

Zasada selekcji $g_0^\mu = 2{,}00$ (vs $2{,}08$)
  & \statuslabel{Program}
  & granica sektora vs $B_{\rm tail}=0$; patrz \texttt{ex60}~T3--T5
  & O-J2 \\

\bottomrule
\end{tabular}
\end{center}

\paragraph{Spójność z duchem teorii.}
Wszystkie wyniki powyżej opierają się wyłącznie na:
\begin{itemize}[nosep]
  \item[$N0$-2] $\alpha = 2$ (z~$K_{ij} = J(\varphi_i\varphi_j)^2$,
    tw.~\ref{prop:substrate-action}),
  \item[$N0$-5] $\beta = \gamma$ (warunek próżni $U'(1)=0$,
    stw.~\ref{prop:vacuum-condition}) $\Rightarrow V(g) = g^3/3 - g^4/4$,
  \item[$N0$-6] $m_{\rm sp} = \sqrt{\gamma}$ (masa bozonu przestrzennego,
    stw.~\ref{prop:N0-6-from-N0-5}).
\end{itemize}
Żaden nowy parametr nie został wprowadzony.
Solitony są rozwiązaniami ODE wynikającej z~akcji TGP --- nie są
zewnętrznymi obiektami wstawionymi do teorii.

\paragraph{Programy obliczeniowe.}

\begin{enumerate}[nosep,label=\textbf{ex\arabic*.}]
  \setcounter{enumi}{59}
  \item \texttt{ex60\_golden\_ratio\_topology.py}: topologia sektorów,
    zera $B_{\rm tail}$, złota proporcja.
  \item \texttt{ex61\_tau\_prediction.py}: predykcja~$g_0^\tau$,
    porównanie kandydatów C1--C5.
  \item \texttt{ex62\_phi\_fixed\_point.py}: skan kinetycznego $\alpha_K$,
    wyznaczenie $\alpha^*_1, \alpha^*_2$ ze warunku
    $F(\alpha_K) = r_{21}$ (szczegóły w Dod.~\ref{app:formula-sumy}).
\end{enumerate}
